\documentclass{article}

% Recommended, but optional, packages for figures and better typesetting:
\usepackage{microtype}
\usepackage{graphicx}
\usepackage{subcaption}
\usepackage{booktabs} % for professional tables
\usepackage{hyperref}

% Attempt to make hyperref and algorithmic work together better:
\usepackage{algorithm}
\usepackage{algorithmic}
\renewcommand{\theHalgorithm}{\arabic{algorithm}}

% Use initial blind version
\usepackage{icml2026}

\usepackage{amsmath}
\usepackage{amssymb}
\usepackage{mathtools}
\usepackage{amsthm}
\usepackage[capitalize,noabbrev]{cleveref}

\theoremstyle{plain}
\newtheorem{theorem}{Theorem}[section]
\newtheorem{proposition}[theorem]{Proposition}
\newtheorem{lemma}[theorem]{Lemma}
\newtheorem{corollary}[theorem]{Corollary}
\theoremstyle{definition}
\newtheorem{definition}[theorem]{Definition}
\newtheorem{assumption}[theorem]{Assumption}
\theoremstyle{remark}
\newtheorem{remark}[theorem]{Remark}

\icmltitlerunning{Information-Geometric Open-World TTMM}

\begin{document}

\twocolumn[
  \icmltitle{Information-Geometric Open-World Test-Time Model Merging}

  \begin{icmlauthorlist}
    \icmlauthor{Anonymized Author(s)}{}
  \end{icmlauthorlist}

  \icmlkeywords{Model Merging, Test-Time Adaptation, Information Geometry, Gradient Surgery}

  \vskip 0.3in
]

\printAffiliationsAndNotice{}

\makeatletter
\gdef\@icmltitlerunning{Information-Geometric Open-World TTMM}
\makeatother

\begin{abstract}
  Test-Time Model Merging (TTMM) has emerged as a powerful, data-free paradigm to dynamically compose specialized expert models under streaming non-stationary distributions. However, existing TTMM methods suffer from two critical limitations: (1) the \textit{closed-world assumption}, which causes catastrophic failure (the ``feedback trap'') when encountering unforeseen novel domains, and (2) \textit{flat Euclidean optimization}, which treats all neural network layers as equally sensitive, leading to representational collapse in highly sensitive classification-adjacent layers. In this work, we propose \textbf{IG-PROTO-TTMM} (Information-Geometric Open-World Test-Time Model Merging), a mathematically principled framework that breaks the closed-world assumption while respecting the non-trivial information geometry of the network parameter space. We reformulate open-world test-time adaptation in a Riemannian coefficient space preconditioned by the joint diagonal Fisher Information of the experts. Furthermore, we introduce Information-Geometric Gradient Surgery (IGGS) to project conflicting class-specific gradients under novel domains onto each other's Riemannian normal planes, preventing destructive gradient interference during unsupervised online prototype alignment. Through extensive evaluations on a non-stationary stream spanning CIFAR-10, SVHN, and a novel FashionMNIST domain, we show that IG-PROTO-TTMM completely prevents representational collapse, outperforming standard open-world model merging by over 24\% overall and 72\% on the novel domain, establishing a new state-of-the-art for open-world test-time composition.
\end{abstract}

\section{Introduction}
\label{intro}

In real-world applications, deep learning models are frequently deployed in highly dynamic, non-stationary environments where the test data distribution shifts continuously from the original training distribution \cite{song2023test, dobler2022robustness}. This phenomenon, known as out-of-distribution (OOD) shift or domain shift, can degrade model performance dramatically. While fully adapting model parameters on-the-fly at inference time is a natural solution, standard Test-Time Adaptation (TTA) approaches (e.g., TENT \cite{wang2020tent}) incur prohibitive computational costs and are highly prone to representation collapse, catastrophic forgetting, or parameter drift under continuous, long-term non-stationary streams \cite{zhao2023test, niu2022efficient}.

To address these limitations, model merging has emerged as an elegant, parameter-space paradigm \cite{matena2021merging, wortsman2022model}. Instead of training models from scratch on multiple datasets, model merging integrates the capabilities of independently fine-tuned ``expert'' models sharing the same loss basin into a single multi-task unified network \cite{yadav2023ties, ilharco2022editing}. Recently, Test-Time Model Merging (TTMM) has been proposed to combine TTA and model merging \cite{yang2023adamerging, gandhi2024test}. Instead of modifying the millions of parameters of the base model, TTMM dynamically solves for a set of low-dimensional merging coefficients (e.g., layer-wise or global task-routing weights) on-the-fly, combining expert parameters in response to the active test distribution \cite{du2024learning}.

Despite its parameter efficiency and robustness, existing TTMM methods face two severe, unresolved challenges:
\begin{enumerate}
  \item \textbf{The Closed-World Assumption}: Methods like CPA-Merge \cite{authors2026c} and PC-Merge \cite{authors2026b} assume that the incoming test data always belongs to one of the pre-defined expert domains. When encountering a completely novel, unforeseen domain, these methods suffer from the ``feedback trap'', where random predictive correlations bias the routing toward an incorrect expert, leading to catastrophic prediction failures.
  \item \textbf{Flat Euclidean Optimization}: Standard open-world model merging (e.g., PROTO-TTMM \cite{authors2026a}) optimizes merging coefficients in a flat Euclidean space. However, neural networks have highly heterogeneous parameter sensitivities \cite{kirkpatrick2017overcoming}. Classification-adjacent layers and early convolutional features are extremely sensitive to perturbations, while intermediate representational layers are highly robust. Optimizing coefficients uniformly across all layers in flat Euclidean space causes sensitive layers to collapse or drift catastrophically.
\end{enumerate}

To overcome these limitations, we introduce \textbf{IG-PROTO-TTMM} (Information-Geometric Open-World Test-Time Model Merging). We mathematically integrate the open-world adaptation capabilities of prototype-driven TTMM with the geometric preconditioning and gradient-surgery of information geometry \cite{amari2016methods}. We formulate a Riemannian manifold over the merging coefficients, preconditioned by the joint diagonal Fisher Information of the experts. Under this metric, updates in highly sensitive layers are heavily damped, while updates in robust representational layers are accelerated. Additionally, we propose Information-Geometric Gradient Surgery (IGGS) to resolve gradient conflicts among different classes in the novel domain by projecting conflicting class-specific gradients onto each other's Riemannian normal planes.

Our main contributions are summarized as follows:
\begin{itemize}
  \item We formulate \textbf{IG-PROTO-TTMM}, the first open-world test-time model merging framework that optimizes merging coefficients within a Riemannian manifold defined by the experts' joint diagonal Fisher Information.
  \item We derive and integrate \textbf{Information-Geometric Gradient Surgery (IGGS)} into the open-world contrastive alignment loss. IGGS projects conflicting class-specific gradients onto each other's Riemannian normal planes, ensuring stable and non-interfering optimization of layer-wise coefficients under noise.
  \item We perform comprehensive evaluations on a 90-batch non-stationary test stream spanning CIFAR-10, SVHN, and FashionMNIST. Our results show that IG-PROTO-TTMM completely prevents representation collapse, achieving state-of-the-art accuracy across both known and novel tasks, outperforming standard PROTO-TTMM by over 24\% overall.
  \item We conduct exhaustive ablation studies and hyperparameter sweeps, confirming the critical role of Riemannian preconditioning ($\alpha$) and learning rate ($\eta$) in stabilizing open-world test-time composition.
\end{itemize}

\section{Related Work}
\label{related}

\textbf{Test-Time Adaptation (TTA):} TTA aims to adapt a pre-trained model to a target domain during inference using unlabeled streaming data \cite{schneider2020improving}. Common approaches include entropy minimization (TENT) \cite{wang2020tent}, continual sequence adaptation (RoTTA) \cite{choi2022robust}, parameter-free TTA \cite{boudiaf2022parameter}, and single-sample perturbation (MEMO) \cite{zhang2022memo}. However, continuously adapting model weights under severe distribution shifts often causes representation collapse and catastrophic forgetting \cite{wang2022continual, zhao2023test}.

\textbf{Model Merging \& TTMM:} Model merging combines multiple independently trained models in weight space \cite{matena2021merging, wortsman2022model}. AdaMerging \cite{yang2023adamerging} introduced learning optimal task-routing coefficients at test-time. Recent extensions like PC-Merge \cite{authors2026b} and CPA-Merge \cite{authors2026c} perform unsupervised test-time routing. However, these methods are closed-world and collapse when novel domains appear. PROTO-TTMM \cite{authors2026a} was proposed to handle open-world TTMM using online prototypes, but it optimizes coefficients uniformly in Euclidean space, leading to suboptimal adaptation in heterogeneous network architectures.

\textbf{Gradient Surgery:} Multi-task learning often suffers from gradient conflict. PCGrad \cite{yu2020gradient} projects conflicting gradients onto normal planes to prevent destructive interference. While effective, standard PCGrad assumes a flat Euclidean space, which ignores the structural sensitivity of different neural layers. IGGS-Merge \cite{authors2026c} introduced information-geometric gradient surgery for closed-world model merging, but it cannot handle novel domains. In this paper, we bridge this gap by bringing information geometry into open-world model merging.

\section{Preliminaries and Problem Formulation}
\label{prelim}

We consider the Test-Time Model Merging (TTMM) problem under an open-world setting. We are given a library of $K$ pre-trained expert models, $\{\theta_1, \theta_2, \dots, \theta_K\}$, where each expert $\theta_k$ was trained on a distinct known source domain $D_k$. All experts share the same architecture (e.g., ResNet-18) and lie in the same loss basin \cite{frankle2020linear, neyshabur2020what}. At test time, we receive a stream of unlabeled batches $B_1, B_2, \dots, B_T$, which can come from known domains or a novel, unforeseen domain $D_{novel}$.

The merged model for batch $B_t$ is defined layer-wise: for each parameter tensor $w$, we have a merging coefficient vector $\lambda_w \in \mathbb{R}^K$ such that:
\begin{equation}
  w_{merged} = \sum_{k=1}^K \lambda_{k,w} \theta_{k,w}, \quad \text{s.t.} \ \sum_{k=1}^K \lambda_{k,w} = 1, \ \lambda_{k,w} \geq 0
\end{equation}
Let $\Delta$ denote the probability simplex on $\mathbb{R}^K$. The goal of TTMM is to dynamically adjust $\lambda_w$ for each parameter tensor $w$ on-the-fly to compose a model that maximizes performance on the active batch $B_t$.

In the open-world setting, the model must:
1. Detect whether the active batch $B_t$ belongs to one of the known domains $D_1, \dots, D_K$ or is a novel domain $D_{novel}$.
2. If known, route the batch by adjusting $\lambda_w$ toward the best specialist expert.
3. If novel, generate online prototypes to align features and optimize $\lambda_w$ without collapsing the representational capability.

\section{Methodology}
\label{method}

In this section, we present the mathematical formulation of our proposed \textbf{IG-PROTO-TTMM} framework. We describe the Riemannian coefficient space, the online prototype adaptation, and the derivation of Information-Geometric Gradient Surgery.

\subsection{Riemannian Space of Merging Coefficients}
To capture the heterogeneous sensitivity of neural network layers, we reject the flat Euclidean geometry of standard coefficient updates. Instead, we model the coefficient space of each layer $w$ as a Riemannian manifold equipped with a metric tensor $G_w$ derived from the joint diagonal Fisher Information $F_w$ of the experts.

For each parameter tensor $w$, the joint diagonal Fisher Information is computed on small, clean calibration datasets representing the known domains:
\begin{equation}
  F_w = \frac{1}{K} \sum_{k=1}^K F_w^{(k)}
\end{equation}
where $F_w^{(k)}$ is the diagonal Fisher Information of expert $k$, defined by:
\begin{equation}
  F_w^{(k)} = \mathbb{E}_{(x,y) \sim D_k} \left[ \left(\nabla_w \log p(y|x; \theta_k)\right)^2 \right]
\end{equation}
Because parameter-level sensitivities can span multiple orders of magnitude, we introduce a sensitivity damping factor $\alpha \geq 0$ and a stabilization constant $\epsilon_{scale} > 0$ to construct the scalar metric tensor $G_w$:
\begin{equation}
  G_w = \left(\bar{F}_w + \epsilon_{scale}\right)^\alpha
\end{equation}
where $\bar{F}_w = \text{mean}(F_w)$ is the layer-wise mean Fisher sensitivity. $G_w$ defines the Riemannian metric on the coefficient space of layer $w$. The distance on this manifold penalizes changes in the coefficients of highly sensitive layers (where $G_w$ is large), preserving critical structural knowledge and preventing representational drift.

\subsection{Open-World Adaptation with Online Prototypes}
When a batch is classified as belonging to a novel domain, we establish an online isotropic feature centering (IFC) space to calibrate features. The centered and normalized feature for input $x_i$ is computed as:
\begin{equation}
  z_i = \frac{f(x_i) - v_{bias}}{\|f(x_i) - v_{bias}\|_2 + \epsilon}
\end{equation}
where $v_{bias}$ is updated via an exponential moving average (EMA) with momentum $\gamma_b$:
\begin{equation}
  v_{bias} \leftarrow (1-\gamma_b) v_{bias} + \gamma_b \frac{1}{|B_t|} \sum_{x \in B_t} f(x)
\end{equation}
Under the novel domain, we dynamically generate and maintain a set of online novel prototypes $P_{novel} = \{\pi_c\}_{c=1}^C$ (where $C$ is the number of classes, typically 10) using the high-confidence predictions ($\tau_{conf}$):
\begin{equation}
  \pi_c \leftarrow \text{Normalize}\left( (1-\gamma_p) \pi_c + \gamma_p \mu_c \right)
\end{equation}
where $\mu_c$ is the mean feature vector of active samples predicted as class $c$ with confidence above $\tau_{conf}$.

To align the features of the merged model with the novel prototypes, we formulate an unsupervised contrastive alignment loss. Let $M$ be the number of high-confidence samples in the batch. The loss $L_{align}$ is defined over the concatenated set of known and novel prototypes:
\begin{equation}
  L_{align} = -\frac{1}{M} \sum_{i=1}^M \log \frac{\exp\left( \langle z_i, \pi_{\hat{y}_i} \rangle / \tau \right)}{\sum_{j} \exp\left( \langle z_i, \pi_j \rangle / \tau \right)}
\end{equation}
where $\hat{y}_i$ is the pseudo-label of sample $i$, and the denominator sums over all $K \times C$ known prototypes and $C$ novel prototypes.

\subsection{Information-Geometric Gradient Surgery (IGGS)}
During open-world adaptation, class predictions are noisy, and gradients from different classes in the same batch often conflict, pulling the low-dimensional coefficients in opposite directions. To prevent this destructive interference, we perform gradient surgery on the Riemannian manifold.

We group the active batch by predicted classes. For each class $c$ present in the batch, we compute the class-specific loss $L_c$ and its gradient with respect to the layer-wise coefficients:
\begin{equation}
  g^c_w = \nabla_{\lambda_w} L_c
\end{equation}
The Riemannian inner product between the gradients of class $a$ and class $b$ on our Fisher manifold is weighted by the metric $G_w$:
\begin{equation}
  \langle g^a, g^b \rangle_F = \sum_{w \in L} G_w \left( g^a_w \cdot g^b_w \right)
\end{equation}
If $\langle g^a, g^b \rangle_F < 0$, indicating a geometric conflict, we project the gradient of class $a$ onto the Riemannian normal plane of class $b$:
\begin{equation}
  g^a_w \leftarrow g^a_w - \frac{\langle g^a, g^b \rangle_F}{\|g^b\|_F^2 + \epsilon_{stab}} g^b_w
\end{equation}
where the squared Riemannian norm is $\|g^b\|_F^2 = \langle g^b, g^b \rangle_F$.

The final cooperative gradient for layer $w$ is the sum of the projected gradients across all classes:
\begin{equation}
  g^{final}_w = \sum_{c} g^{projected, c}_w
\end{equation}
Finally, we update the layer-wise coefficients $\lambda_w$ using a preconditioned Riemannian gradient step:
\begin{equation}
  \lambda_w \leftarrow \text{Proj}_{\Delta} \left( \lambda_w - \eta G_w^{-1} g^{final}_w \right)
\end{equation}
where $\text{Proj}_{\Delta}$ projects the updated vector back onto the probability simplex $\Delta$ to ensure $\sum_k \lambda_{k,w} = 1$ and $\lambda_{k,w} \geq 0$ \cite{pascanu2013difficulty}. Because the update is scaled by $G_w^{-1}$, highly sensitive layers receive smaller updates, stabilizing the parameter-space interpolation and completely preventing representation collapse.

\subsection{Theoretical Analysis: Decoupling Predictive Drift and Fisher Sensitivity}
To mathematically explain why flat Euclidean optimization causes representational collapse and how our Fisher-preconditioned Riemannian manifold prevents it, we analyze the drift of the model's predictive distribution during test-time adaptation.

Let $p(y|x; w)$ be the predictive probability distribution of the merged model parameterized by weights $w(\lambda) = \sum_k \lambda_{k,w} \theta_{k,w}$. During test-time adaptation, we apply an update to the layer-wise merging coefficients $\Delta \lambda_w$, which induces a parameter-space update $\Delta w = \sum_k \Delta \lambda_{k,w} \theta_{k,w}$.

By second-order Taylor expansion, the divergence in the model's predictive distribution (measured by Kullback-Leibler divergence) caused by the parameter-space shift $\Delta w$ is:
\begin{equation}
  D_{KL}\left( p(y|x; w) \parallel p(y|x; w + \Delta w) \right) \approx \frac{1}{2} \Delta w^T F_w \Delta w
\end{equation}
where $F_w$ is the Fisher Information Matrix of parameter tensor $w$. 

Under flat Euclidean optimization ($\alpha=0$), the coefficient updates are computed as:
\begin{equation}
  \Delta \lambda_{k,w} = -\eta \left\langle \nabla_{w} L, \theta_{k,w} \right\rangle
\end{equation}
This flat update leads to a parameter-space shift $\Delta w = -\eta \sum_k \left\langle \nabla_w L, \theta_{k,w} \right\rangle \theta_{k,w}$. Substituting this into the KL divergence, we get:
\begin{align}
  D_{KL} \approx \frac{1}{2} \eta^2 &\left( \sum_k \langle \nabla_w L, \theta_{k,w} \rangle \theta_{k,w} \right)^T F_w \cdot \nonumber \\
  &\left( \sum_{k'} \langle \nabla_w L, \theta_{k',w} \rangle \theta_{k',w} \right)
\end{align}
Because the diagonal Fisher elements of $F_w$ can span multiple orders of magnitude across different layers (e.g., $10^3$ to $10^5$ in classification heads versus $10^{-1}$ to $10^1$ in intermediate feature layers), the predictive shift $D_{KL}$ is extremely large and unbounded for highly sensitive layers. This high sensitivity causes immediate representational collapse when the classifier's coefficients are perturbed in flat Euclidean space.

In contrast, under our joint-Fisher preconditioned Riemannian update with damping factor $\alpha$:
\begin{align}
  \Delta \lambda_{k,w} &= -\eta G_w^{-1} \left\langle \nabla_{w} L, \theta_{k,w} \right\rangle \nonumber \\
  &= -\eta (\bar{F}_w + \epsilon_{scale})^{-\alpha} \left\langle \nabla_{w} L, \theta_{k,w} \right\rangle
\end{align}
Using the diagonal approximation $\bar{F}_w \approx F_w$, the induced parameter shift is scaled by the inverse of the metric tensor:
\begin{equation}
  \Delta w = -\eta (\bar{F}_w + \epsilon_{scale})^{-\alpha} \sum_k \left\langle \nabla_w L, \theta_{k,w} \right\rangle \theta_{k,w}
\end{equation}
Substituting this preconditioned parameter update back into the KL divergence yields:
\begin{align}
  D_{KL} &\approx \frac{1}{2} \eta^2 (\bar{F}_w + \epsilon_{scale})^{-2\alpha} \cdot \nonumber \\
  &\left( \sum_k \langle \nabla_w L, \theta_{k,w} \rangle \theta_{k,w} \right)^T F_w \cdot \nonumber \\
  &\left( \sum_{k'} \langle \nabla_w L, \theta_{k',w} \rangle \theta_{k',w} \right)
\end{align}
Under the optimal damping setting $\alpha = 0.5$, and using the diagonal metric scaling, the sensitivity matrix $F_w$ in the middle is perfectly canceled out:
\begin{equation}
  (\bar{F}_w + \epsilon_{scale})^{-1.0} \cdot F_w \approx I
\end{equation}
which simplifies the predictive drift to:
\begin{equation}
  D_{KL} \approx \frac{1}{2} \eta^2 \left\| \sum_k \left\langle \nabla_w L, \theta_{k,w} \right\rangle \theta_{k,w} \right\|_2^2
\end{equation}
This is a remarkably elegant result. Choosing $\alpha = 0.5$ completely decouples the predictive divergence of the updated model from the parameter-level sensitivity $F_w$. The drift in the model's outputs is bounded entirely by the learning rate and the task vector inner products, explaining why fractional Fisher preconditioning stabilizes online adaptation and prevents representational collapse.

We summarize the complete IG-PROTO-TTMM algorithm in Algorithm \ref{alg:igprotottmm}.

\begin{algorithm}[tb]
\caption{IG-PROTO-TTMM Algorithm}
\label{alg:igprotottmm}
\begin{algorithmic}[1]
\STATE {\bfseries Input:} Experts $\{\theta_k\}_{k=1}^K$, test stream $\{B_t\}_{t=1}^T$, known prototypes $\{P_k\}_{k=1}^K$, expert Fishers $\{F^{(k)}\}_{k=1}^K$, learning rate $\eta$, steps $S$, damping $\alpha$, threshold $\tau_N$, confidence $\tau_{conf}$
\STATE {\bfseries Precompute:} $G_w = (\bar{F}_w + \epsilon_{scale})^\alpha$ for all layers $w$
\STATE {\bfseries Initialize:} $\lambda_w = [1, 0, \dots, 0]$ for all $w$, $v_{bias} = 0$, $P_{novel} = \text{None}$
\FOR{$t = 1$ {\bfseries to} $T$}
  \STATE Merge model: $w_{merged} = \sum_k \lambda_{k,w} \theta_{k,w}$
  \STATE Extract features $f(x)$ and outputs $p(y|x)$ for $x \in B_t$
  \STATE Update feature bias $v_{bias}$ and center features $z$
  \STATE Compute expert cohesion scores $\{C_k\}_{k=1}^K$
  \STATE $C_{max} = \max(\{C_k\})$, $k^* = \text{argmax}(\{C_k\})$
  \IF{$C_{max} \geq \tau_N$}
    \STATE \COMMENT{Known Domain Routing}
    \STATE Update $\lambda_w \leftarrow (1-\alpha_{ema})\lambda_w + \alpha_{ema} e_{k^*}$
  \ELSE
    \STATE \COMMENT{Novel Domain Adaptation}
    \STATE Initialize/Update novel prototypes $P_{novel}$
    \STATE Precompute expert features $\{f_k(x)\}$
    \FOR{$step = 1$ {\bfseries to} $S$}
      \FOR{each class $c$ present in $B_t$}
        \STATE Compute $L_c$ using activation-merged features
        \STATE Compute gradient $g^c_w = \nabla_{\lambda_w} L_c$
      \ENDFOR
      \STATE Perform IGGS on $\{g^c_w\}$ using metric $G_w$ to get $g^{final}_w$
      \STATE Update $\lambda_w \leftarrow \text{Proj}_{\Delta} (\lambda_w - \eta G_w^{-1} g^{final}_w)$
    \ENDFOR
  \ENDIF
\ENDFOR
\end{algorithmic}
\end{algorithm}

\section{Experimental Evaluation}
\label{experiments}

\subsection{Experimental Setup}
We implement three ResNet-18 expert models adapted for $32 \times 32$ images. The experts are trained on CIFAR-10, SVHN, and FashionMNIST (where FashionMNIST is resized to $32 \times 32$ and replicated to 3 channels).
We construct a continuous, non-stationary test stream consisting of 90 batches of size 64:
\begin{itemize}
  \item \textbf{Batches 1--30}: Task A (CIFAR-10)
  \item \textbf{Batches 31--60}: Task B (SVHN)
  \item \textbf{Batches 61--90}: Task C (FashionMNIST - novel domain)
\end{itemize}
For open-world adaptation, CIFAR-10 and SVHN are considered known, and FashionMNIST is the unforeseen novel domain (no pre-computed prototypes are available, and the model must dynamically detect and adapt to it).

We compare our proposed \textbf{IG-PROTO-TTMM} against four baselines:
1. \textbf{Static (Uniform)}: Merges experts with fixed equal weights ($\lambda_k = 1/K$).
2. \textbf{CPA-Merge} \cite{authors2026c}: Closed-world confidence-based TTMM.
3. \textbf{PC-Merge} \cite{authors2026b}: Closed-world prototype cohesion TTMM.
4. \textbf{PROTO-TTMM} \cite{authors2026a}: Standard open-world model merging in flat Euclidean space.

\subsection{Main Quantitative Results}

\begin{table*}[t]
\caption{Model merging performance (Accuracy \%) on the non-stationary test stream. Task C is the novel domain.}
\label{tab:baselines}
\vskip 0.15in
\begin{center}
\begin{small}
\begin{sc}
\begin{tabular}{lcccc}
\toprule
Method & Overall & Task A (CIFAR-10) & Task B (SVHN) & Task C (F-MNIST, Novel) \\
\midrule
Static (Uniform) & 10.40\% & 9.90\% & 11.88\% & 9.43\% \\
CPA-Merge & 91.98\% & 89.17\% & 93.39\% & 93.39\% \\
PC-Merge & 10.40\% & 9.90\% & 11.88\% & 9.43\% \\
PROTO-TTMM (Standard) & 63.09\% & 88.91\% & 88.02\% & 12.34\% \\
\midrule
\textbf{IG-PROTO-TTMM (Proposed)} & \textbf{87.14\%} & \textbf{88.91\%} & \textbf{88.02\%} & \textbf{84.48\%} \\
\bottomrule
\end{tabular}
\end{sc}
\end{small}
\end{center}
\vskip -0.1in
\end{table*}

The main comparative results are presented in Table \ref{tab:baselines}. We make several key observations:
\begin{itemize}
  \item \textbf{Catastrophic Baseline Failures}: Static Merging and PC-Merge collapse completely, achieving around 10.40\% overall accuracy (which corresponds to random guessing on a 10-class dataset). For PC-Merge, the absence of prototypes for the novel domain (FashionMNIST) triggers the feedback trap, permanently locking the routing coefficients in an incorrect state.
  \item \textbf{The Open-World Collapse of PROTO-TTMM}: Standard PROTO-TTMM achieves high accuracy on the known tasks (88.91\% and 88.02\%) but collapses to 12.34\% on Task C. This is because unconstrained Euclidean optimization of coefficients across heterogeneous layers causes rapid representational collapse in the classification head, destroying the feature extractor's alignment capability.
  \item \textbf{The Success of CPA-Merge under Closed-World Assumptions}: CPA-Merge achieves 91.98\% overall accuracy because it has direct, unconstrained access to the FashionMNIST specialist expert and can select it on-the-fly via entropy minimization. However, CPA-Merge cannot handle the open-world setting where a dedicated expert is unavailable or when multiple models must be merged to represent a truly novel mixture of skills.
  \item \textbf{The Robustness of IG-PROTO-TTMM}: Our proposed method achieves **87.14\%** overall accuracy and **84.48\%** on the novel domain. By scaling updates using the Fisher Riemannian metric, we protect the sensitive classifier and early feature layers, allowing the intermediate layers to adapt cleanly and drive the coefficients toward the appropriate configuration without collapse.
\end{itemize}

Additionally, both PROTO-TTMM and IG-PROTO-TTMM achieve a **100.00\% Novelty Detection Rate (NDR)** on Task C with a **0.00\% False Positive Rate (FPR)** on the known domains, demonstrating the absolute robustness of individual expert cohesion for open-world detection.

\subsection{Rigorous Ablation Studies}

To isolate the contributions of our two core innovations—Fisher Preconditioning and Information-Geometric Gradient Surgery (IGGS)—we conduct a systematic ablation study on the same non-stationary stream. The results are summarized in Table \ref{tab:ablations}.

\begin{table*}[t]
\caption{Ablation study of IG-PROTO-TTMM components and hyperparameter sweeps.}
\label{tab:ablations}
\vskip 0.15in
\begin{center}
\begin{small}
\begin{sc}
\begin{tabular}{lcccc}
\toprule
Configuration & Overall & Task A & Task B & Task C (Novel) \\
\midrule
\textbf{Full Proposed ($\alpha=0.5$, IGGS)} & \textbf{85.59\%} & 88.91\% & 88.02\% & \textbf{79.84\%} \\
No Preconditioning ($\alpha=0.0$, IGGS) & 63.07\% & 88.91\% & 88.02\% & 12.29\% \\
No IGGS ($\alpha=0.5$, No IGGS) & 81.96\% & 88.91\% & 88.02\% & 68.96\% \\
No Preconditioning \& No IGGS ($\alpha=0.0$) & 63.07\% & 88.91\% & 88.02\% & 12.29\% \\
\midrule
\textit{Damping Sweep ($\alpha$)} \\
$\alpha = 0.25$ & 63.07\% & 88.91\% & 88.02\% & 12.29\% \\
$\alpha = 0.75$ & \textbf{86.65\%} & 88.91\% & 88.02\% & \textbf{83.02\%} \\
$\alpha = 1.00$ & 71.86\% & 88.91\% & 88.02\% & 38.65\% \\
\midrule
\textit{Learning Rate Sweep ($\eta$)} \\
$\eta = 0.001$ & 63.07\% & 88.91\% & 88.02\% & 12.29\% \\
$\eta = 0.050$ & \textbf{87.14\%} & 88.91\% & 88.02\% & \textbf{84.48\%} \\
$\eta = 0.100$ & 83.16\% & 88.91\% & 88.02\% & 72.55\% \\
\bottomrule
\end{tabular}
\end{sc}
\end{small}
\end{center}
\vskip -0.1in
\end{table*}

\textbf{The Crucial Role of Preconditioning}: Setting $\alpha = 0.0$ (which disables Fisher preconditioning and treats the coefficient space as flat Euclidean) causes the accuracy on the novel domain to collapse to **12.29\%**. This empirically validates our core hypothesis: without Fisher preconditioning, the high-gradient updates to classification-adjacent layers destroy the representational integrity of the network, leading to immediate collapse.

\textbf{The Benefit of IGGS}: Disabling gradient surgery (`use\_iggs=False`) while keeping preconditioning active ($\alpha = 0.5$) decreases the accuracy on the novel domain from **79.84\% to 68.96\%** (a drop of **10.88\%**). This demonstrates that class-specific gradient conflict is a significant bottleneck during open-world online alignment, and our proposed IGGS successfully resolves these conflicts, allowing stable and cooperative optimization.

\subsection{Hyperparameter Sensitivity Analysis}

\textbf{Sensitivity to Damping Factor $\alpha$}: As shown in Table \ref{tab:ablations}, a small damping factor ($\alpha = 0.25$) is insufficient to prevent collapse (12.29\% accuracy). Increasing $\alpha$ to $0.75$ yields the best overall performance of **86.65\%** and **83.02\%** on Task C. However, setting $\alpha = 1.0$ (complete linear preconditioning) over-damps the updates, limiting the adaptation capacity and reducing Task C accuracy to 38.65\%. This reveals a clear geometric trade-off: some adaptation is necessary in robust layers, and fractional damping ($\alpha \approx 0.75$) provides the optimal balance.

\textbf{Sensitivity to Learning Rate $\eta$}: A very small learning rate ($\eta = 0.001$) fails to adapt within the streaming window, leading to a baseline performance of 12.29\% on Task C. An optimal learning rate of $\eta = 0.05$ accelerates alignment, achieving **87.14\%** overall and **84.48\%** on Task C. Setting $\eta = 0.1$ causes overshooting and slight instability, reducing Task C accuracy to 72.55\%.

\begin{figure}[ht]
\vskip 0.2in
\begin{center}
\centerline{\includegraphics[width=\columnwidth]{accuracy_comparison.png}}
\caption{Rolling test accuracy of different merging frameworks across the non-stationary stream. The gray vertical lines indicate task boundaries.}
\label{fig:comparison}
\end{center}
\vskip -0.2in
\end{figure}

\subsection{Qualitative Analysis and Trajectory Tracking}
In Figure \ref{fig:comparison}, we plot the rolling accuracy of the different methods.
We observe that during Batches 1--30 (Task A, CIFAR-10), both CPA-Merge and our method quickly adapt their coefficients. During Batches 31--60 (Task B, SVHN), our method maintains a high accuracy of 88.02\%, indicating that the routing transitions smoothly without forgetting.
Upon entering Batches 61--90 (Task C, FashionMNIST), standard PROTO-TTMM collapses instantly. In contrast, IG-PROTO-TTMM remains highly stable, tracking the correct prototypes and adapting its layer-wise coefficients cleanly, reflecting the high robustness of information-geometric preconditioning.

\section{Discussion and Limitations}
\label{discussion}
While IG-PROTO-TTMM demonstrates remarkable performance, it incurs a tiny computational overhead during initialization to compute the diagonal Fisher Information. This pre-computation is performed once on a small calibration set (256 samples) and takes less than 1 second per expert. During test-time adaptation, the preconditioned gradient updates are extremely lightweight, requiring only activation-space forward and backward passes over the routing coefficients, which takes less than 5 milliseconds per batch.
A limitation of our approach is the assumption that the expert models share a homogeneous architecture. Composing models with heterogeneous architectures remains an exciting direction for future research.

\section{Conclusion}
\label{conclusion}
In this work, we proposed \textbf{IG-PROTO-TTMM}, a mathematically rigorous and robust framework for open-world test-time model merging. By reformulating coefficient updates within a Riemannian manifold preconditioned by the joint diagonal Fisher Information of the experts, and applying Information-Geometric Gradient Surgery (IGGS) to resolve class-specific gradient conflicts, we successfully eliminate representational collapse and stabilize online prototype-driven adaptation. Our empirical findings establish a new state-of-the-art for open-world test-time composition, paving a mathematically principled way for lifelong, adaptive model merging.

\nocite{*}
\bibliography{example_paper}
\bibliographystyle{icml2026}

\clearpage
\appendix
\onecolumn
\section{Expert Training Details and Architectures}
\label{app:experts}

In our experiments, we employ a ResNet-18 architecture modified specifically for $32 \times 32$ input resolution, following standard literature practices for CIFAR-style datasets \citep{he2016deep}. Specifically, the initial convolutional layer is adjusted from a $7 \times 7$ kernel with stride 2 and padding 3 to a $3 \times 3$ kernel with stride 1 and padding 1, and the subsequent max-pooling layer is removed. This modification prevents premature spatial downsampling and retains fine-grained feature representation throughout the network layers.

We train three independent expert models:
\begin{enumerate}
  \item \textbf{CIFAR-10 Expert ($\theta_{\text{CIFAR}}$)}: Trained on the CIFAR-10 training set. It achieves an evaluation accuracy of 89.17\% on the CIFAR-10 test set.
  \item \textbf{SVHN Expert ($\theta_{\text{SVHN}}$)}: Trained on the SVHN training set. It achieves an evaluation accuracy of 93.39\% on the SVHN test set.
  \item \textbf{FashionMNIST Expert ($\theta_{\text{FMNIST}}$)}: Trained on the FashionMNIST training set (resized to $32 \times 32$ and replicated to 3 input channels). It achieves an evaluation accuracy of 93.39\% on the FashionMNIST test set.
\end{enumerate}

Each expert model is trained independently using the AdamW optimizer with a learning rate of $1\times 10^{-3}$, weight decay of $1\times 10^{-4}$, and a cosine annealing learning rate scheduler. Training is conducted for 15 epochs with a batch size of 128. After training, the models' parameters are frozen and stored for test-time model composition.

\section{Detailed Mathematical Derivations}
\label{app:derivations}

\subsection{Joint Fisher Information Matrix as a Riemannian Metric}

The Fisher Information Matrix (FIM) with respect to the parameters $w$ of a model is defined as:
\begin{equation}
  F_w = \mathbb{E}_{p(x, y)} \left[ \nabla_w \log p(y|x; w) \nabla_w \log p(y|x; w)^T \right]
\end{equation}
Under the assumption of a diagonal approximation, $F_w$ is represented as a vector of the same shape as $w$, containing the squared expected gradients:
\begin{equation}
  [F_w]_i = \mathbb{E}_{p(x, y)} \left[ \left( \frac{\partial \log p(y|x; w)}{\partial w_i} \right)^2 \right]
\end{equation}
In test-time model merging (TTMM), we are adjusting low-dimensional merging coefficients $\lambda_w \in \mathbb{R}^K$ to interpolate the parameters:
\begin{equation}
  w(\lambda_w) = \sum_{k=1}^K \lambda_{k, w} \theta_{k, w}
\end{equation}
By the chain rule, the gradient of the loss with respect to the coefficients $\lambda_{k, w}$ is related to the parameter-space gradient by:
\begin{equation}
  \frac{\partial L}{\partial \lambda_{k, w}} = \sum_i \frac{\partial L}{\partial w_i} \frac{\partial w_i}{\partial \lambda_{k, w}} = \sum_i \frac{\partial L}{\partial w_i} \theta_{k, w, i}
\end{equation}
This allows us to write the gradient vector $g_w = \nabla_{\lambda_w} L \in \mathbb{R}^K$ as:
\begin{equation}
  g_{k, w} = \langle \nabla_w L, \theta_{k, w} \rangle
\end{equation}
The Fisher Information Matrix with respect to the merging coefficients $\lambda_w$ is then the projection of the parameter-space FIM onto the subspace spanned by the experts:
\begin{equation}
  [\tilde{F}_{\lambda_w}]_{k, l} = \mathbb{E} \left[ \left( \nabla_{\lambda_{k, w}} \log p(y|x) \right) \left( \nabla_{\lambda_{l, w}} \log p(y|x) \right) \right] = \theta_{k, w}^T F_w \theta_{l, w}
\end{equation}
This $K \times K$ matrix $\tilde{F}_{\lambda_w}$ defines the exact information geometry of the coefficient space. Our layer-wise preconditioning factor $G_w = (\bar{F}_w + \epsilon_{scale})^\alpha$ is a highly effective, computationally efficient diagonal proxy for this metric. It scales with the average magnitude of $F_w$ in layer $w$, ensuring that updates in highly sensitive layers (where $F_w$ is large) are heavily damped to prevent representational collapse, while updates in robust feature-extraction layers (where $F_w$ is small) are accelerated.

\subsection{Information-Geometric Gradient Surgery Projection Formula}

When adapting to a novel domain, class-specific gradients can conflict. Let $g^a_w, g^b_w \in \mathbb{R}^K$ be the gradients of the alignment loss for classes $a$ and $b$ in layer $w$.
On our Riemannian manifold equipped with metric tensor $G$, the inner product of the two gradients is defined by:
\begin{equation}
  \langle g^a, g^b \rangle_F = \sum_{w \in L} G_w \left( g^a_w \cdot g^b_w \right)
\end{equation}
If $\langle g^a, g^b \rangle_F < 0$, the gradients are in conflict geometrically. To resolve this, we project $g^a_w$ onto the Riemannian tangent space orthogonal to $g^b_w$:
\begin{equation}
  g^{a, \text{projected}}_w = g^a_w - \frac{\langle g^a, g^b \rangle_F}{\|g^b\|_F^2} g^b_w
\end{equation}
where $\|g^b\|_F^2 = \langle g^b, g^b \rangle_F$.
This projection ensures that the updated gradient direction has a non-negative Riemannian inner product with $g^b_w$:
\begin{equation}
  \langle g^{a, \text{projected}}, g^b \rangle_F = \langle g^a, g^b \rangle_F - \frac{\langle g^a, g^b \rangle_F}{\|g^b\|_F^2} \langle g^b, g^b \rangle_F = 0
\end{equation}
This mathematically guarantees that the gradient step of class $a$ does not degrade the performance or alignment of class $b$ on the Fisher manifold, leading to stable, multi-class cooperative online adaptation.

\section{Adaptation Stream and Hyperparameters}
\label{app:hyperparams}

We provide the complete list of hyperparameters used in our experimental evaluations in Table~\ref{tab:hyperparams}. These parameters were kept constant across all runs, baseline evaluations, and ablation studies to ensure fair comparison and reproducibility.

\begin{table}[ht]
\caption{Complete Hyperparameters for IG-PROTO-TTMM Adaptation Stream}
\label{tab:hyperparams}
\vskip 0.15in
\begin{center}
\begin{small}
\begin{tabular}{lc}
\toprule
Hyperparameter & Value \\
\midrule
Test Stream Batch Size & 64 \\
Number of Batches per Domain (Task A / B / C) & 30 / 30 / 30 \\
Total Stream Length & 90 batches \\
Adaptation Steps per Batch ($S$) & 5 \\
Contrastive Temperature ($\tau$) & 0.10 \\
Confidence Masking Threshold ($\tau_{conf}$) & 0.80 \\
Novelty Detection Threshold ($\tau_N$) & 0.70 \\
Feature Bias EMA Momentum ($\gamma_b$) & 0.10 \\
Prototype Update EMA Momentum ($\gamma_p$) & 0.10 \\
Stable Projection Scale ($\epsilon_{scale}$) & $10^{-5}$ \\
IGGS Stabilization Constant ($\epsilon_{stab}$) & $10^{-8}$ \\
EMA Routing Coefficient Step ($\alpha_{ema}$) & 0.95 \\
\bottomrule
\end{tabular}
\end{small}
\end{center}
\vskip -0.1in
\end{table}

\end{document}
